\chapter{Materiais e Métodos}\label{cap3-materiais-metodos}

Este capítulo descreve o ambiente computacional, o \emph{framework} experimental e o protocolo de execução utilizados para comparar as quatro variantes do algoritmo memético AG+\(2\)-opt apresentadas na \secref{sec-14-objetivos-especificos}.
O objetivo é permitir que outros pesquisadores reproduzam, com o máximo de precisão possível, os experimentos discutidos no \capref{cap4-resultados}, respeitando as limitações de hardware disponíveis.

\section{Ambiente computacional e \emph{framework} experimental}\label{sec-31-ambiente-framework}

Esta seção detalha as características do hardware e do software utilizados, bem como a arquitetura do \emph{framework} experimental desenvolvido para garantir a integridade e a reprodutibilidade das comparações.

\subsection{Hardware e restrições físicas}\label{sec-311-hardware-software}

Os experimentos foram conduzidos em um \emph{notebook} equipado com processador \textbf{Intel Core i7-7700HQ} e uma GPU \textbf{NVIDIA GeForce GTX 1050 Mobile}.
Trata-se de uma GPU dedicada de entrada para \emph{notebooks}, baseada na arquitetura \textbf{Pascal}.

Embora o uso deste equipamento tenha sido motivado inicialmente pela disponibilidade, o conjunto experimental oferece uma vantagem metodológica: a combinação de uma CPU \emph{quad-core} relativamente mais capaz com uma GPU de porte modesto acentua o custo de comunicação e os limites de memória, tornando as diferenças entre estratégias mais nítidas.
Por outro lado, a seleção experimental foi limitada a instâncias de até \(n_{coords} = \maiorinstancia\) cidades (\texttt{pr1002}) devido ao tempo de execução da variante de referência em CPU.
Como o protocolo estatístico exige \(n_{reps} = \totalrepeticoes\) execuções independentes de cada variante por problema, instâncias acima desse porte tornariam proibitivo o tempo de coleta da implementação puramente em CPU, cujo processamento sequencial dos indivíduos na busca local \(2\)-opt --- somado ao crescimento do número de gerações até a convergência --- eleva acentuadamente o custo computacional total.

As especificações relevantes do ambiente computacional utilizado nos experimentos estão sintetizadas na \tabref{tbl:hardware-specs}, destacando o contraste entre os \(4\) núcleos (\(8\) \emph{threads}) da CPU e os \(640\) núcleos CUDA da GPU distribuídos em \(5\) SMs (\emph{Streaming Multiprocessors}).

\begin{table}[htb]
      \centering
      \caption{Especificações do ambiente computacional utilizado nos experimentos}\label{tbl:hardware-specs}
      \begin{tabelapadrao}
            \begin{tabular}{@{}
                        >{\raggedright\arraybackslash}p{(\linewidth - 4\tabcolsep) * \real{0.2600}}
                        >{\raggedright\arraybackslash}p{(\linewidth - 4\tabcolsep) * \real{0.3400}}
                        >{\raggedright\arraybackslash}p{(\linewidth - 4\tabcolsep) * \real{0.4000}}@{}}
                  \toprule
                  Característica        & CPU (Intel Core i7-7700HQ)                                              & GPU (NVIDIA GeForce GTX 1050 Mobile)                     \\
                  \midrule
                  Arquitetura / Geração & Kaby Lake (2017)                                                        & Pascal GP107 (2016, Compute Capability 6.1)              \\
                  Núcleos / Unidades    & \(4\) núcleos / \(8\) \emph{threads}                                    & \(640\) núcleos CUDA em \(5\) SMs                        \\
                  Frequência            & \(2{,}80\)--\(3{,}80\text{ GHz}\)                                       & \(1354\)--\(1493\text{ MHz}\)                            \\
                  Memória               & \(16\text{ GB}\) RAM DDR4                                               & \(4\text{ GB}\) VRAM GDDR5 (\(\approx 112\text{ GB/s}\)) \\
                  Interface CPU--GPU    & \multicolumn{2}{c@{}}{Barramento PCIe 3.0 (\(\approx 16\text{ GB/s}\))}                                                            \\
                  \bottomrule
            \end{tabular}
      \end{tabelapadrao}
\end{table}

O desempenho computacional é estritamente limitado pelas fronteiras físicas da GPU.
Detalha-se a seguir o impacto dessas restrições no processamento:

\begin{enumerate}
      \item
            \textbf{Memória de vídeo (VRAM) e memória compartilhada:} A GPU dispõe de \(4\text{ GB}\) de VRAM.
            Seguindo a convenção da TSPLIB \autocite{reinelt1991tsplib}, tanto a matriz completa de distâncias \(D\) (\(n_{coords}^2\) elementos) quanto os dois \emph{buffers} de população (\(2 \times n_{pop} \times n_{coords}\) elementos, com \(n_{pop} = 256\)) são armazenados como inteiros de \(32\) bits (\texttt{int32}, \(4\) bytes cada), enquanto o vetor de custos usa inteiros de \(64\) bits (\(256 \times 8 = 2048\) bytes).
            Na variante Totalmente em GPU, que mantém todo o estado residente na placa, o consumo de VRAM é dado por:
            \begin{equation}\label{eq:vram-total}
                  Mem_{total}(n_{coords}) \approx 4\,n_{coords}^2 + 2048\,n_{coords} + 2048 \text{ bytes.}
            \end{equation}
            Mesmo na maior instância testada (\(n_{coords} = \maiorinstancia\), \texttt{pr1002}), a Equação~\eqref{eq:vram-total} resulta em cerca de \(6{,}07\text{ MB}\) --- uma fração mínima dos \(4\text{ GB}\) disponíveis.
            Além disso, a janela de memória compartilhada dos blocos comporta rotas de até \(2048\) cidades, de modo que todas as \totalinstancias\ instâncias deste estudo cabem inteiramente na memória rápida local sem necessidade de particionamento.
      \item
            \textbf{Comunicação CPU--GPU via barramento PCIe:} A troca de dados entre CPU e GPU ocorre pelo barramento PCIe, cuja largura de banda (\(\approx 16\text{ GB/s}\)) é cerca de sete vezes menor que a largura de banda interna da VRAM (\(\approx 112\text{ GB/s}\)).
            Transferências frequentes penalizam diretamente o tempo total de execução conforme a decomposição:
            \begin{equation}\label{eq:tempo-total}
                  T_{total} = T_{cpu} + T_{gpu} + T_{transf}\text{,}
            \end{equation}
            podendo anular o ganho obtido nos núcleos gráficos.
            Esse contraste motiva reduzir o número de transferências por geração (\textbf{Híbrida Otimizada}) ou manter todo o ciclo evolutivo na GPU (\textbf{Totalmente em GPU}), conforme detalhado na \secref{sec-322-variantes}.
      \item
            \textbf{Ocupação dos \(640\) núcleos CUDA (\(5\) SMs):} Manter centenas de \emph{threads} ativas simultaneamente nos \(5\) \emph{SMs}\footnote{\textbf{Limite arquitetural da GPU utilizada:} Na microarquitetura Pascal (Compute Capability 6.1) da GTX~1050 Mobile, cada SM suporta no máximo \(2048\) \emph{threads} residentes (\(64\) \emph{warps}); em outras gerações de GPU esse teto difere.
            O uso elevado de registradores ou de memória compartilhada por bloco reduz a ocupação efetiva abaixo desse limite.}
            (\emph{Streaming Multiprocessors}) --- cada um com 128 cores, totalizando 640 núcleos de execução --- permite que os multiprocessadores alternem a execução entre grupos de \emph{threads} prontos enquanto outros aguardam leituras de memória, ocultando a latência de acesso (\emph{latency hiding}) e sustentando alta taxa de ocupação dos \(640\) núcleos \autocite{kirk2010programming}.
\end{enumerate}

\subsection{Ambiente de software}\label{sec-312-ambiente-software}

O \emph{framework} experimental foi desenvolvido em \texttt{Python 3.12.13} sobre \texttt{Debian GNU/Linux 13}, utilizando um conjunto de bibliotecas e compiladores selecionados para assegurar consistência funcional entre as execuções em CPU e GPU.
As versões exatas, os requisitos mínimos de arquitetura CUDA e o papel de cada componente central constam na \tabref{tbl:software-specs}.

\begin{table}[htb]
      \centering
      \caption{Especificações, requisitos mínimos e papel dos componentes do ambiente de software}\label{tbl:software-specs}
      \begin{tabelapadrao}
            \begin{tabular}{@{}
                        >{\raggedright\arraybackslash}p{(\linewidth - 4\tabcolsep) * \real{0.2600}}
                        >{\raggedright\arraybackslash}p{(\linewidth - 4\tabcolsep) * \real{0.3400}}
                        >{\raggedright\arraybackslash}p{(\linewidth - 4\tabcolsep) * \real{0.4000}}@{}}
                  \toprule
                  Componente                     & Versão / Requisito                                             & Papel no ambiente                                                     \\
                  \midrule
                  Sistema operacional            & Debian GNU/Linux 13 (\emph{trixie}), kernel 6.12.43 (\texttt{x86\_64}) & Plataforma de execução do \emph{benchmark}                            \\
                  \sepitem
                  \emph{Driver} gráfico NVIDIA   & 580.196.07                                                     & Interface com a GPU e suporte ao CUDA Runtime 12.9                    \\
                  \sepitem
                  CUDA Toolkit / compilador      & CUDA 12.9 (\texttt{nvcc} V12.9.86)                             & Compilação e execução dos \emph{kernels} CUDA customizados            \\
                  \sepitem
                  Compute Capability CUDA        & Mínimo \(\ge 5{,}0\) (testado: \(6{,}1\), Pascal)              & Requisito mínimo do CUDA 12.9 e do \texttt{cupy-cuda12x}              \\
                  \sepitem
                  Compilador C/C++ (\emph{host}) & GCC 12.4.0                                                     & Compilador auxiliar de \emph{host} exigido pelo \texttt{nvcc}         \\
                  \sepitem
                  Interpretador Python           & Python 3.12.13 (mínimo \(\ge 3{,}12\))                         & Orquestração dos experimentos e interface com CPU/GPU                 \\
                  \sepitem
                  Computação numérica            & NumPy 2.5.4; \hyperlink{term:cupy}{CuPy} 14.2.0                & \emph{Backend} de CPU (NumPy) e GPU (CuPy sobre CUDA 12.9)            \\
                  \sepitem
                  Estatística e banco de dados   & SciPy 1.18.1; scikit-posthocs 0.17.1; DuckDB 2.0.0             & Protocolo de \textcite{demsar2006statistical} e consolidação (\apref{apendice-banco}) \\
                  \bottomrule
            \end{tabular}
      \end{tabelapadrao}
\end{table}

A compatibilidade de API entre NumPy e CuPy foi determinante para manter a lógica algorítmica idêntica entre as variantes de CPU e GPU, alterando apenas o \emph{backend} de execução.
Os demais pacotes auxiliares de análise, visualização e validação declarados no projeto, bem como o procedimento reproduzível de instalação do ambiente via \texttt{uv}, estão documentados no \apref{apendice-reproducao}.

Embora a API de alto nível do CuPy seja adequada para operações matriciais regulares, na variante \textbf{Totalmente em GPU} (\texttt{FullGPUModular}) tanto a busca local \(2\)-opt quanto os operadores genéticos (seleção por torneio, cruzamento OX, mutação por troca e avaliação de custo) são executados por \textbf{kernels CUDA customizados} escritos em C++ (\texttt{RawKernel}).
O desenvolvimento direto desses \emph{kernels} evita sobrecargas repetidas de chamadas da API de alto nível e viabiliza três ganhos estruturais:

\begin{enumerate}
      \item
            \textbf{Fusão e encadeamento de etapas na GPU:} Executar o refinamento \(2\)-opt, a avaliação de custos e os operadores evolutivos diretamente na GPU, evitando retornos intermediários à CPU.
      \item
            \textbf{Uso explícito da memória rápida local:} Armazenar dados ativos das rotas e estruturas auxiliares na memória compartilhada dos blocos e em registradores, reduzindo acessos à memória global.
      \item
            \textbf{Cooperação entre \emph{threads} na busca \(2\)-opt:} Em vez de limitar cada rota a uma única \emph{thread} isolada, os \emph{kernels} customizados permitem que múltiplas \emph{threads} de um mesmo bloco cooperem na avaliação dos movimentos candidatos de cada rota e processem toda a matriz de rotas da população de uma só vez.
\end{enumerate}

A avaliação de custos usa \textbf{aritmética inteira}.
O arredondamento da TSPLIB é aplicado uma única vez, ao montar a matriz de distâncias \(D\), por \(\text{nint}(x) = \lfloor x + 0{,}5 \rfloor\), conforme o tipo de cada instância (por exemplo, a distância euclidiana em \texttt{EUC\_2D}, com regras próprias em \texttt{CEIL\_2D}, \texttt{ATT} e \texttt{GEO}) \autocite{reinelt1991tsplib}.
A partir desse ponto, todas as distâncias entre cidades são números inteiros de \(32\) bits (\texttt{int32}) e o custo total de qualquer circuito é acumulado em \(64\) bits (\texttt{int64}), o que evita ultrapassar o limite de representação de \(32\) bits em rotas de grande porte.
Consequentemente, o custo de qualquer rota candidata e o ganho \(\Delta C\) de qualquer troca \(2\)-opt reduzem-se a somas e subtrações exatas de inteiros, eliminando divergências de arredondamento de ponto flutuante entre CPU e GPU e garantindo que duas variantes quaisquer atribuam \textbf{exatamente o mesmo custo} a uma mesma rota.

\subsection{Arquitetura do \emph{framework} experimental}\label{sec-313-arquitetura-framework}

Para garantir comparações justas, o código foi estruturado em uma arquitetura modular que separa a lógica dos algoritmos da orquestração dos experimentos.
O \emph{framework} opera em três camadas distintas, ilustradas na \figref{fig:class-diagram}:

\begin{enumerate}
      \item
            \textbf{Camada de Algoritmos:} Adota um padrão modular de estratégias intercambiáveis em torno de uma única classe orquestradora (\texttt{GeneticAlgorithm}), cujas seis etapas do ciclo evolutivo recebem implementações de CPU ou de GPU sob o mesmo contrato; a definição detalhada das quatro variantes é apresentada na \secref{sec-32-arquitetura-variantes}.
      \item
            \textbf{Camada de Orquestração (\hyperlink{term:benchmark}{Benchmark}):} Responsável por carregar as configurações experimentais (parâmetros, instâncias), gerenciar a execução das repetições independentes e garantir o isolamento entre testes.
            Esta camada implementa mecanismos de \emph{checkpoint} para salvar o estado de cada execução, permitindo a recuperação em caso de falhas e a auditoria posterior dos dados.
      \item
            \textbf{Camada de Análise Estatística e Visualização:} Scripts dedicados a processar os \emph{checkpoints}, consolidar os dados e gerar as tabelas, diagramas e relatórios estatísticos.
\end{enumerate}

\begin{figure}[htbp]
      \centering
      \caption{Arquitetura modular em três camadas do \emph{framework} experimental}\label{fig:class-diagram}
      \begin{tikzpicture}[node distance=5mm and 4mm]
            % CAMADA 1: Algoritmos e Estrategias
            \node[fc_proc, text width=12.6cm] (solver) {\textbf{GeneticAlgorithm}\\(\texttt{+evolve()}, \texttt{+timer: PhaseTimer}, composição de estratégias)};
            \node[fc_grouptitle, above=2.5mm of solver] (tc1) {Camada 1 — Algoritmos e Estratégias Intercambiáveis};
            \node[fc_proc, below=5.5mm of solver, text width=3.7cm] (s2) {\texttt{Avaliação}\\(\emph{EvaluationStrategy})};
            \node[fc_proc, left=4mm of s2, text width=3.7cm] (s1) {\texttt{Inicialização}\\(\emph{InitializationStrategy})};
            \node[fc_proc, right=4mm of s2, text width=3.7cm] (s3) {\texttt{Seleção}\\(\emph{SelectionStrategy})};
            \node[fc_proc, below=3.5mm of s1, text width=3.7cm] (s4) {\texttt{Cruzamento OX}\\(\emph{CrossoverStrategy})};
            \node[fc_proc, below=3.5mm of s2, text width=3.7cm] (s5) {\texttt{Mutação Swap}\\(\emph{MutationStrategy})};
            \node[fc_memetic, below=3.5mm of s3, text width=3.7cm] (s6) {\texttt{Busca Local 2-opt}\\(\emph{LocalSearchStrategy})};
            \node[fc_term, below=9mm of s5, text width=12.6cm] (vars) {\textbf{4 Variantes Instanciadas:}\\CPU \quad\textbar\quad Híbrida Ingênua \quad\textbar\quad Híbrida Otimizada \quad\textbar\quad Totalmente em GPU};

            % CAMADA 2: Execução e Armazenamento (titulo centralizado dentro do topo, seta central sem colisão)
            \node[fc_proc, below=24mm of vars, text width=3.7cm, minimum height=14mm] (l2b) {\textbf{Pré-cálculo}\\Matriz \(D\) (\texttt{int32}) e\\vizinhos (\(k_{\text{viz}} = \kviz\))};
            \node[fc_proc, left=4mm of l2b, text width=3.7cm, minimum height=14mm] (l2a) {\textbf{Carregadores}\\Leitura \texttt{TSPLIB} e\\consulta \texttt{DuckDB}};
            \node[fc_proc, right=4mm of l2b, text width=3.7cm, minimum height=14mm] (l2c) {\textbf{Orquestração}\\\(n_{reps} = \totalrepeticoes\) sementes,\\JSON \(\to\) \texttt{DuckDB}};
            \node[fc_grouptitle, above=2.5mm of l2b] (tc2) {Camada 2 — Execução e Armazenamento};

            % CAMADA 3: Analise Estatística e Visualização (simétrico em x=0: (5.9+0.5)/2 = 3.2cm)
            \node[fc_proc, below=24mm of l2b, xshift=-3.2cm, text width=5.9cm, minimum height=14mm] (l3a) {\textbf{Protocolo Demšar (2006)}\\Friedman \(F_F\), pós-teste\\de Nemenyi (\(CD\)) e Wilcoxon\\pareado + Holm--Bonferroni};
            \node[fc_proc, right=5mm of l3a, text width=5.9cm, minimum height=14mm] (l3b) {\textbf{Geração Automática}\\Tabelas \LaTeX,\\diagramas \(CD\) em TikZ\\e gráficos vetoriais PDF/SVG};
            \node[fc_grouptitle, above=2.5mm of $(l3a.north)!0.5!(l3b.north)$] (tc3) {Camada 3 — Análise Estatística e Visualização};

            % Caixas de camada
            \begin{scope}[on background layer]
                  \node[draw=black!35, dashed, rounded corners=2pt, fill=white, inner sep=4pt, fit=(s1)(s2)(s3)(s4)(s5)(s6)] (sgrid) {};
                  \node[fc_group, fit=(tc1)(solver)(sgrid)(vars)] (c1) {};
                  \node[fc_group, fit=(tc2)(l2a)(l2b)(l2c)] (c2) {};
                  \node[fc_group, fit=(tc3)(l3a)(l3b)] (c3) {};
            \end{scope}

            % Redesenha títulos por cima do fundo
            \node[fc_grouptitle] at (tc1) {Camada 1 — Algoritmos e Estratégias Intercambiáveis};
            \node[fc_grouptitle] at (tc2) {Camada 2 — Execução e Armazenamento};
            \node[fc_grouptitle] at (tc3) {Camada 3 — Análise Estatística e Visualização};

            % Conexões centralizadas e estritamente verticais
            \draw[fc_arr] (solver.south) -- (sgrid.north);
            \draw[fc_arr] (sgrid.south) -- (vars.north);
            \draw[fc_arr] (c1.south) -- node[fc_lbl, right=2mm, pos=0.48] {Configuração e execução das variantes} (c1.south |- c2.north);
            \draw[fc_arr] (c2.south) -- node[fc_lbl, right=2mm, pos=0.48] {Checkpoints JSON / registros DuckDB} (c2.south |- c3.north);
      \end{tikzpicture}
\end{figure}

Esta separação assegura que as métricas de tempo e qualidade sejam coletadas de forma consistente para todas as variantes, eliminando vieses que poderiam surgir de implementações \emph{ad hoc} para cada plataforma.

\section{Variantes algorítmicas e estratégias de paralelização}\label{sec-32-arquitetura-variantes}

As variantes comparadas neste trabalho são \textbf{análogas}.
Todas partem do mesmo algoritmo memético, com os mesmos operadores, os mesmos hiperparâmetros globais e os mesmos critérios de parada, e diferem em \emph{onde} cada etapa é executada --- CPU ou GPU --- e em \emph{como} os dados trafegam e são processados: sequencialmente na CPU, por circuito individual \(C_m\) ou pela matriz completa da população \(P(t)\) (\secref{sec-322-variantes}).
Paralelizar muda o modo de execução, e é por isso que as variantes são tratadas como análogas, e não como equivalentes.
A assimetria de avaliação entre CPU e GPU --- bem como a equivalência bit a bit verificada entre as variantes híbridas --- é detalhada na \secref{sec-5311-semantica-operadores}.

\subsection{O algoritmo genético base (AG + \texorpdfstring{\(2\)-opt}{2-opt})}\label{sec-321-ga-base}

O algoritmo base é um Algoritmo Genético (AG) com seleção elitista de sobreviventes do tipo \((\mu + \lambda)\), na qual pais e descendentes têm o mesmo tamanho (\(\mu = \lambda = n_{pop}\)), hibridizado com uma busca local \(2\)-opt truncada.
Esta combinação, denominada Algoritmo Memético, equilibra a exploração global do espaço de busca (via operadores genéticos) com a intensificação local (via \(2\)-opt).
Todos os hiperparâmetros mantidos constantes nas quatro variantes constam na \tabref{tbl:ga-hyperparameters} e são detalhados a seguir:

\begin{table}[htbp]
      \centering
      \caption{Hiperparâmetros e configurações fixas do algoritmo memético AG+\(2\)-opt adotados nas quatro variantes}\label{tbl:ga-hyperparameters}
      \begin{tabelapadrao}
            \begin{tabular}{@{}
                        >{\raggedright\arraybackslash}p{(\linewidth - 4\tabcolsep) * \real{0.3400}}
                        >{\centering\arraybackslash}p{(\linewidth - 4\tabcolsep) * \real{0.1600}}
                        >{\raggedright\arraybackslash}p{(\linewidth - 4\tabcolsep) * \real{0.5000}}@{}}
                  \toprule
                  Parâmetro / Componente                  & Símbolo                    & Valor / Especificação                                                                      \\
                  \midrule
                  Tamanho da população                    & \(n_{pop}\)                & \(256\) rotas                                                                              \\
                  Seleção de progenitores                 & \(\tau_{\text{trnmt}}\)    & Torneio de tamanho \(5\)                                                                   \\
                  Cruzamento (\emph{Order Crossover}, OX) & \(p_c\)                    & \(1{,}0\) (aplicado a todos os \(n_{pop}/2\) pares)                                        \\
                  Mutação por troca (\emph{swap})         & \(p_{\text{mut}}\)         & \(0{,}02\) por rota descendente                                                            \\
                  Busca local \(2\)-opt (pós-mutação)     & \(\kappa\)                 & Até \(\maxpassadas\) varreduras \emph{best-improvement}                                    \\
                  Lista estática de candidatos            & \(k_{\text{viz}}\)         & \(\kviz\) vizinhos mais próximos por cidade                                                \\
                  Seleção de sobreviventes (elitismo)     & \((\mu + \lambda)\)        & Determinística com \(\mu = \lambda = 256\) (melhores \(256\) de \(512\))                   \\
                  Critério de parada (i) --- alvo         & \(\varepsilon_{\text{alvo}}\) & \(\text{gap} \le 1\%\) em relação a \(B^*\)                                             \\
                  Critério de parada (ii) --- estagnação  & \(G_{\text{estag}}\)       & \(\lceil 2\sqrt{n_{coords}} \rceil\) gerações sem melhoria                                 \\
                  Critério de parada (iii) --- teto       & \(G_{\max}\)               & \(\lceil 2\,n_{coords}\sqrt{n_{coords}} \rceil\) gerações                                  \\
                  Repetições e sementes                   & \(n_{reps}\)               & \(\totalrepeticoes\) execuções (sementes inteiras \(\sementeinicial\) a \(\sementefinal\)) \\
                  \bottomrule
            \end{tabular}
      \end{tabelapadrao}
\end{table}

As escolhas de projeto subjacentes aos parâmetros da \tabref{tbl:ga-hyperparameters} obedecem a três diretrizes técnicas:

\begin{itemize}
      \item
            \textbf{Escala de memória da população (\(n_{pop} = 256\)):} O tamanho fixo enquadra-se na faixa usual de \(50\) a \(500\) indivíduos da literatura de AGs para o TSP \autocite{harik1999gambler,goldberg1992massive,larranaga1999genetic} e equilibra a ocupação em VRAM. Enquanto a matriz de distâncias cresce quadraticamente (\(\mathcal{O}(n_{coords}^2)\), ocupando \(4\,n_{coords}^2\) bytes), a população cresce bilinearmente (\(\mathcal{O}(n_{pop}\,n_{coords})\), ocupando \(4\,n_{pop}\,n_{coords}\) bytes por \emph{buffer}), mantendo o custo de uma população abaixo de \(1\text{ MiB}\) (\(\approx 1{,}03\text{ MB}\)) até \(n_{coords} = \maiorinstancia\) e evitando que um \(n_{pop}\) proporcional a \(n_{coords}\) introduza um segundo termo quadrático na memória do dispositivo.
      \item
            \textbf{Representação e operadores evolutivos:} As rotas são codificadas como permutações de inteiros sem repetição; a seleção por torneio independe de reescalonamento global de aptidão e paraleliza-se diretamente por sorteios independentes; o operador OX preserva a ordem relativa de sub-rotas válidas ao recombinar cada par de progenitores; e a mutação por troca (\emph{swap}) introduz perturbações mínimas mantendo a factibilidade do circuito.
      \item
            \textbf{Acoplamento pós-mutação e lista estática de vizinhos (\(k_{\text{viz}} = \kviz\)):} A busca local \(2\)-opt \emph{best-improvement} é aplicada sobre cada descendente pós-mutação \autocite{alsheddy2017memetic,larranaga1999genetic} com vizinhança restrita aos \(k_{\text{viz}} = \kviz\) vizinhos mais próximos (\secref{sec-221-2opt}), encerrando-se o ciclo evolutivo assim que qualquer um dos três critérios de parada listados na \tabref{tbl:ga-hyperparameters} é satisfeito.
\end{itemize}

Para justificar empiricamente a escolha de \(k_{\text{viz}} = \kviz\), realizou-se um microbenchmark isolado com \(64\) rotas e \(\kappa = \maxpassadas\) varreduras de \(2\)-opt comparando a lista de candidatos contra a varredura exaustiva completa, cujos resultados constam na \tabref{tbl:candidate-k-microbenchmark}.
Em instâncias muito pequenas (\texttt{berlin52}), o custo de indireção da tabela empata com a varredura direta (\(0{,}99\times\) na CPU; \(0{,}79\times\) na GPU), mas já a partir de \(n_{coords} = 100\) a restrição de vizinhança acelera a busca local em até \(8{,}22\times\) na CPU (\texttt{pr264}) e \(4{,}05\times\) na GPU (\texttt{lin318}), viabilizando a execução completa da linha de base em CPU até \(n_{coords} = \maiorinstancia\) sob exatamente o mesmo orçamento de busca.

\begin{table}[htbp]
      \centering
      \caption{Microbenchmark isolado da restrição por lista de candidatos (\(k_{\text{viz}} = \kviz\) vs.\ varredura exaustiva, \(64\) rotas, \(\kappa = \maxpassadas\) passadas)}\label{tbl:candidate-k-microbenchmark}
      \begin{tabelapadrao}
            \begin{tabular}{@{}lrrrrrr@{}}
                  \toprule
                  Instância                    & \multicolumn{3}{c}{CPU (s)} & \multicolumn{3}{c}{GPU -- matriz de rotas (s)}                                                                                  \\
                  \cmidrule(lr){2-4} \cmidrule(lr){5-7}
                                               & \(k_{\text{viz}} = \kviz\)  & Exaustiva                                      & Ganho            & \(k_{\text{viz}} = \kviz\) & Exaustiva   & Ganho            \\
                  \midrule
                  \texttt{berlin52} (\(n=52\)) & \(0{,}935\)                 & \(0{,}924\)                                    & \(0{,}99\times\) & \(0{,}001\)                & \(0{,}001\) & \(0{,}79\times\) \\
                  \texttt{kroA100} (\(n=100\)) & \(1{,}590\)                 & \(4{,}250\)                                    & \(2{,}67\times\) & \(0{,}002\)                & \(0{,}003\) & \(1{,}51\times\) \\
                  \texttt{ch150} (\(n=150\))   & \(3{,}122\)                 & \(8{,}066\)                                    & \(2{,}58\times\) & \(0{,}003\)                & \(0{,}005\) & \(2{,}03\times\) \\
                  \texttt{pr264} (\(n=264\))   & \(4{,}808\)                 & \(39{,}513\)                                   & \(8{,}22\times\) & \(0{,}004\)                & \(0{,}016\) & \(3{,}72\times\) \\
                  \texttt{lin318} (\(n=318\))  & \(5{,}811\)                 & \(45{,}798\)                                   & \(7{,}88\times\) & \(0{,}006\)                & \(0{,}023\) & \(4{,}05\times\) \\
                  \bottomrule
            \end{tabular}
      \end{tabelapadrao}
\end{table}

\subsection{Estratégias de paralelização (as quatro variantes)}\label{sec-322-variantes}

Para investigar o impacto da GPU, o algoritmo base foi instanciado em quatro variantes, que não são implementações separadas: todas usam o mesmo algoritmo, e a variante define apenas quais implementações ocupam cada uma das seis estratégias intercambiáveis do ciclo evolutivo --- inicialização, seleção, cruzamento, mutação, busca local e avaliação.
Conectando o ciclo conceitual de sete etapas (\secref{sec-251-ga-2opt}) à decomposição temporal em oito fases cronometradas no \capref{cap4-resultados} (\secref{sec-44-decomposicao-fases} e \apref{apendice-telemetria}), a sétima etapa --- seleção de sobreviventes elitista \((\mu + \lambda)\) --- executa na CPU nas três primeiras variantes e na GPU na variante Totalmente em GPU, enquanto a oitava fase registra explicitamente a sincronização entre CPU e GPU (\texttt{Sinc.}).

\begin{enumerate}
      \item
            \textbf{CPU (\hyperlink{term:baseline}{\emph{baseline}}):} Execução inteiramente na CPU utilizando NumPy: inicialização, seleção, cruzamento, mutação, busca local e avaliação apontam para as implementações NumPy da mesma classe \texttt{GeneticAlgorithm}.
            No nível arquitetural, a CPU permanece sendo a plataforma MIMD de referência; o que ocorre aqui é o aproveitamento pontual de vetorização SIMD dentro dos núcleos para acelerar partes da avaliação dos movimentos \(2\)-opt.
            Já o processamento dos indivíduos, no nível do algoritmo, ocorre essencialmente de forma sequencial.
            Esta variante serve como referência para medir o ganho de performance absoluto.
      \item
            \textbf{HybridNaive (Híbrida Ingênua):} Mantém a população, os operadores genéticos e a avaliação de custos na CPU, enviando apenas a busca local \(2\)-opt para a GPU.
            A cada geração, reenvia a matriz de distâncias \(D\) e transfere cada uma das \(n_{pop} = 256\) rotas individualmente (uma a uma) para a placa, disparando um lançamento de \emph{kernel} por rota e trazendo-a de volta em seguida.
            É esse envio individual --- que multiplica as transferências (\(\mathit{H2D}\) e \(\mathit{D2H}\)) e os lançamentos de \emph{kernel} em cada geração --- que caracteriza a abordagem como ``ingênua'', expondo a latência de disparo e o gargalo do barramento PCIe.
      \item
            \textbf{HybridOptimized (Híbrida Otimizada):} Também mantém os operadores genéticos na CPU, mas fixa a matriz \(D\) em VRAM e agrupa todas as rotas da população \(P(t)\) em uma única matriz (\(n_{pop} \times n_{coords}\)).
            Em cada geração, toda a população é enviada em bloco à GPU e refinada por um único lançamento do \emph{kernel} \(2\)-opt, que devolve a matriz refinada à CPU; em seguida, a avaliação de custos também roda em bloco na GPU, retornando apenas o vetor com os \(256\) custos.
            Comparada à Híbrida Ingênua, essa organização em lote elimina o reenvio repetido de \(D\) e reduz drasticamente o número de lançamentos de \emph{kernel}.
      \item
            \textbf{FullGPUModular (Totalmente em GPU):} Implementa a estratégia de ``GPU residente'': a matriz de distâncias \(D\), a lista de candidatos \(N_k\) e a população inicial \(P(0)\) (gerada sob a mesma semente compartilhada) são enviadas uma única vez na inicialização à VRAM, onde permanecem residentes enquanto todas as etapas iterativas do ciclo evolutivo --- seleção, cruzamento, mutação, busca local \(2\)-opt, avaliação e seleção de sobreviventes \((\mu + \lambda)\) --- executam na GPU.
            Essa organização elimina completamente o retorno de matrizes de rotas para a CPU (\(0{,}00\)\,MB em \(\mathit{D2H}\), restrito a escalares de parada) e reduz o tráfego recorrente \(\mathit{H2D}\) a apenas cerca de \(11{,}26\)\,KB por geração (independentemente de \(n_{coords}\)), correspondentes aos sorteios aleatórios gerados na CPU pelo NumPy (participantes dos torneios, pontos de corte do OX e posições de mutação) para manter exatamente a mesma sequência de sementes nas quatro variantes.
            Os operadores de GPU não usam as abstrações de alto nível do CuPy, e sim \emph{kernels} CUDA escritos especificamente para este trabalho e pré-compilados antes das medições.
\end{enumerate}

\section{Seleção de instâncias e protocolo experimental}\label{sec-33-selecao-instancias}

Para garantir a relevância estatística dos resultados, foram selecionadas \totalinstancias\ instâncias da biblioteca \textbf{TSPLIB} \autocite{reinelt1991tsplib}, variando de \(\menorinstancia\) a \(\maiorinstancia\) cidades, de modo a abranger diferentes escalas de dimensão no conjunto avaliado.
O conjunto inclui problemas com diferentes características geométricas e diferentes formatos de cálculo de distância (\texttt{EUC\_2D}, \texttt{ATT}, \texttt{GEO} e \texttt{EXPLICIT}) \autocite{reinelt1991tsplib}, permitindo avaliar os algoritmos em diferentes cenários.
As \totalinstancias\ instâncias selecionadas estão listadas em ordem crescente de dimensão na \tabref{tbl:selected-instances}, juntamente com seus respectivos formatos de distância e ótimos conhecidos (\(B^*\)).

\begin{tabelapadrao}
\begin{longtable}[]{@{}
            >{\raggedright\arraybackslash}p{(\linewidth - 6\tabcolsep) * \real{0.2400}}
            >{\centering\arraybackslash}p{(\linewidth - 6\tabcolsep) * \real{0.1800}}
            >{\centering\arraybackslash}p{(\linewidth - 6\tabcolsep) * \real{0.2800}}
            >{\raggedleft\arraybackslash}p{(\linewidth - 6\tabcolsep) * \real{0.3000}}@{}}
      \caption{Instâncias da TSPLIB selecionadas para o estudo}\label{tbl:selected-instances}\tabularnewline
      \toprule\noalign{}
      Instância         & \(n_{coords}\) & Tipo de distância & Ótimo conhecido (\(B^*\)) \\
      \midrule\noalign{}
      \endfirsthead
      \toprule\noalign{}
      Instância         & \(n_{coords}\) & Tipo de distância & Ótimo conhecido (\(B^*\)) \\
      \midrule\noalign{}
      \endhead
      \bottomrule\noalign{}
      \endlastfoot
      \texttt{att48}    & \(48\)         & \texttt{ATT}      & \(10628\)                 \\*
      \texttt{eil51}    & \(51\)         & \texttt{EUC\_2D}  & \(426\)                   \\*
      \texttt{berlin52} & \(52\)         & \texttt{EUC\_2D}  & \(7542\)                  \\*
      \texttt{gr96}     & \(96\)         & \texttt{GEO}      & \(55209\)                 \\*
      \texttt{kroA100}  & \(100\)        & \texttt{EUC\_2D}  & \(21282\)                 \\*
      \texttt{ch150}    & \(150\)        & \texttt{EUC\_2D}  & \(6528\)                  \\*
      \texttt{d198}     & \(198\)        & \texttt{EUC\_2D}  & \(15780\)                 \\*
      \texttt{pr264}    & \(264\)        & \texttt{EUC\_2D}  & \(49135\)                 \\*
      \texttt{lin318}   & \(318\)        & \texttt{EUC\_2D}  & \(42029\)                 \\*
      \texttt{pcb442}   & \(442\)        & \texttt{EUC\_2D}  & \(50778\)                 \\*
      \texttt{att532}   & \(532\)        & \texttt{ATT}      & \(27686\)                 \\*
      \texttt{si535}    & \(535\)        & \texttt{EXPLICIT} & \(48450\)                 \\*
      \texttt{gr666}    & \(666\)        & \texttt{GEO}      & \(294358\)                \\*
      \texttt{rat783}   & \(783\)        & \texttt{EUC\_2D}  & \(8806\)                  \\*
      \texttt{pr1002}   & \(1002\)       & \texttt{EUC\_2D}  & \(259045\)                \\
\end{longtable}
\end{tabelapadrao}

O protocolo de execução seguiu as diretrizes para comparação de algoritmos estocásticos:

\begin{itemize}
      \item
            \textbf{Repetições:} Cada par (algoritmo, instância) foi executado em \textbf{\totalrepeticoes\ repetições} independentes.
      \item
            \textbf{Sementes Aleatórias:} As sementes foram fixadas e distintas por repetição --- derivadas da semente base \(\sementeinicial\) somada ao índice da repetição (de \(\sementeinicial\) a \(\sementefinal\)) ---, o que torna cada execução reprodutível e amostra adequadamente o comportamento estocástico.
\end{itemize}

\section{Protocolo de análise estatística}\label{sec-34-protocolo-estatistico}

Aplicando os testes não paramétricos formulados na \secref{sec-27-comparacao-estatistica} \autocite{demsar2006statistical} ao nível de significância \(\alpha = 0{,}05\), avaliam-se as duas métricas centrais do estudo: a qualidade das soluções (\emph{gap} percentual frente a \(B^*\)) e o desempenho computacional (tempo de execução e \hyperlink{term:speedup}{\emph{speedup}}).
O protocolo combina dois eixos complementares sobre as \(m_{\text{alg}} = \totalvariantes\) variantes, as \(n_{\text{inst}} = \totalinstancias\) instâncias e as \(n_{reps} = \totalrepeticoes\) repetições independentes, sintetizados na \figref{fig:results-flowchart}.

\subsection{Eixo 1: comparações pareadas por instância}\label{sec-341-eixo1}

Dentro de cada uma das \(n_{\text{inst}} = \totalinstancias\) instâncias, confrontam-se todos os \(m = 6\) pares possíveis entre as \totalvariantes\ variantes sobre as \(n_{reps} = \totalrepeticoes\) execuções pareadas pelas mesmas sementes aleatórias, utilizando o teste de Wilcoxon da Equação~\eqref{eq:wilcoxon} \autocite{wilcoxon1945individual}.
Para controlar o erro do Tipo~I global nas \(m = 6\) comparações simultâneas de cada instância, os valores-\(p\) recebem o ajuste sequencial de Holm--Bonferroni (\secref{sec-272-comparacoes-multiplas}), comparando-se os valores-\(p\) ordenados aos limiares \(\alpha / (6 - i + 1)\) para \(i = 1, \dots, 6\) \autocite{holm1979simple}.
Quando mais da metade das diferenças pareadas em um par é exatamente nula (empates idênticos), registra-se o número efetivo de diferenças não nulas como alerta de poder estatístico reduzido \autocite[Seção~3.1.3]{demsar2006statistical}.

\subsection{Eixo 2: comparação global multi-instância}\label{sec-342-eixo2}

Na análise agregada, cada combinação algoritmo--instância é resumida pela média das \(n_{reps} = \totalrepeticoes\) repetições, gerando uma matriz \(N \times k\) com \(N = n_{\text{inst}} = \totalinstancias\) e \(k = m_{\text{alg}} = \totalvariantes\) para o \emph{gap} e outra para o tempo de execução.
Sobre cada matriz, calculam-se os postos médios \(R_j\) pela Equação~\eqref{eq:friedman-ranks}, a estatística de Friedman \(\chi_F^2\) pela Equação~\eqref{eq:friedman-chi2} e a estatística \(F_F\) de Iman--Davenport pela Equação~\eqref{eq:iman-davenport} (com \(3\) e \(42\) graus de liberdade) \autocite{friedman1937use,demsar2006statistical}.

Sempre que a hipótese nula global \(H_0\) é rejeitada (\(p_F < 0{,}05\)), aplica-se o pós-teste de Nemenyi pela Equação~\eqref{eq:nemenyi-cd} para comparar todos os \(6\) pares de variantes \autocite{nemenyi1963distribution}.
Substituindo \(k = \totalvariantes\), \(N = \totalinstancias\) e o valor crítico tabelado \(q_{0{,}05} = 2{,}569\) \autocite[Tabela~5]{demsar2006statistical} na Equação~\eqref{eq:nemenyi-cd}, obtém-se o limiar \emph{a priori} \(CD = 1{,}211\), utilizado na \figref{fig:results-flowchart} e nos diagramas de diferença crítica do \capref{cap4-resultados}.
Por fim, ao contrastar exclusivamente as três variantes em GPU frente à referência em CPU ao longo das \totalinstancias\ instâncias (\(m = 3\) hipóteses), emprega-se o teste pareado de Wilcoxon com ajuste de Holm--Bonferroni, conforme recomendado por \textcite[Seção~3.2.2]{demsar2006statistical}.

\begin{figure}[htbp]
      \centering
      \caption{Fluxograma unificado do protocolo experimental e de análise estatística dos resultados}\label{fig:results-flowchart}
      \begin{tikzpicture}[node distance=5.5mm and 8mm]
            % Topo comum
            \node[fc_term, text width=12.6cm] (collect) {\textbf{Coleta Experimental:} \(n_{\text{inst}} = \totalinstancias\) instâncias TSPLIB \(\times\) \(m_{\text{alg}} = \totalvariantes\) variantes \(\times\) \(n_{reps} = \totalrepeticoes\) execuções independentes};
            \node[fc_proc, below=5mm of collect, text width=12.6cm] (metrics) {\textbf{Métricas Extraídas:} Qualidade (\(\text{gap}\) percentual frente a \(B^*\)) e Desempenho (\(T_{\text{total}}\), tempo por fase e \emph{speedup})};

            % Eixo 1 (Esquerda): Analise Pareada
            \node[fc_proc, below left=22mm and -5.0cm of metrics, text width=5.1cm] (p1) {Comparações nos 6 pares 2 a 2\\(\(\totalrepeticoes\) repetições por instância)};
            \node[fc_proc, below=6mm of p1, text width=5.1cm] (p2) {Teste de postos sinalizados de\\Wilcoxon + ajuste Holm--Bonferroni\\(\(\alpha = 0{,}05\))};
            \node[fc_term, below=6mm of p2, text width=5.1cm] (p3) {Significância estatística pareada\\em cada instância};

            % Eixo 2 (Direita): Analise Global
            \node[fc_proc, below right=22mm and -5.9cm of metrics, text width=5.8cm] (g1) {Teste Omnibus de Friedman com\\ estatística \(F_F\) de Iman--Davenport\\(\(\alpha = 0{,}05\), \(n_{\text{inst}} = \totalinstancias\))};
            \node[fc_dec, below=5.5mm of g1, text width=2.7cm] (gdec) {Rejeita \(H_0\)\\(\(p_F < 0{,}05\))?};
            \node[fc_proc, below left=6.5mm and -0.8cm of gdec, text width=2.9cm] (gnem) {Pós-teste de Nemenyi\\(\(CD = 1{,}211\)) e\\Diagramas \(CD\)};
            \node[fc_proc, below right=6.5mm and -0.8cm of gdec, text width=2.9cm] (gno) {Sem diferença\\significativa\\(\(p_F \ge 0{,}05\))};

            % Caixas de fundo + títulos fc_grouptitle centralizados no topo com entrada simétrica deslocada da borda externa
            \begin{scope}[on background layer]
                  \node[fc_group, fit=(p1)(p2)(p3)] (box1) {};
                  \node[fc_group, fit=(g1)(gdec)(gnem)(gno)] (box2) {};
            \end{scope}
            \node[fc_grouptitle, fill=white, anchor=south] at (box1.north) {Eixo 1 — Pareada (\totalrepeticoes\ reps)};
            \node[fc_grouptitle, fill=white, anchor=south] at (box2.north) {Eixo 2 — Global (Demšar, 2006)};

            % Arestas ortogonais simétricas entrando nas bordas externas livres do titulo
            \draw[fc_arr] (collect) -- (metrics);
            \draw[fc_arr] (metrics.south) -- ++(0,-4.5mm) -| ($(box1.north west)+(0.6cm,0)$);
            \draw[fc_arr] (metrics.south) -- ++(0,-4.5mm) -| ($(box2.north east)+(-0.6cm,0)$);
            \draw[fc_arr] (p1) -- (p2);
            \draw[fc_arr] (p2) -- (p3);
            \draw[fc_arr] (g1) -- (gdec);
            \draw[fc_arr] (gdec) -| node[fc_lbl, pos=0.25] {Sim} (gnem);
            \draw[fc_arr] (gdec) -| node[fc_lbl, pos=0.25] {Não} (gno);
      \end{tikzpicture}
\end{figure}
